\documentclass[12pt,a4paper]{article}

% Pacotes de codificação e linguagem
\usepackage[utf8]{inputenc}
\usepackage[T1]{fontenc}
\usepackage[brazil]{babel}
\usepackage{csquotes}

% Pacotes de formatação
\usepackage[margin=2.5cm,top=3cm,bottom=2cm]{geometry}
\usepackage{setspace}
\usepackage{parskip}
\usepackage{titlesec}

% Pacotes gráficos e cores
\usepackage{graphicx}
\usepackage{xcolor}
\usepackage{booktabs}
\usepackage{longtable}
\usepackage{colortbl}

% Pacotes para links e referências
\usepackage[pdftex]{hyperref}
\usepackage{abntex2cite}

% Configuração de cores para links
\hypersetup{
    colorlinks=true,
    linkcolor=blue,
    filecolor=magenta,
    urlcolor=blue,
    citecolor=blue,
    pdftitle={Projeto de Rede Industrial: Fábrica de Tijolos Ecológicos},
    pdfauthor={Felipe Marques},
    pdfsubject={Redes Industriais},
}

% Configuração dos títulos das seções
\titleformat{\section}
  {\normalfont\Large\bfseries}{\thesection}{1em}{}
\titleformat{\subsection}
  {\normalfont\large\bfseries}{\thesubsection}{1em}{}
\titleformat{\subsubsection}
  {\normalfont\normalsize\bfseries}{\thesubsubsection}{1em}{}

% Espaçamento entre linhas
\onehalfspacing


% Início do documento
\begin{document}

% Capa
\begin{titlepage}
\centering
\vspace*{2cm}

{\Huge\bfseries Projeto de Rede Industrial\par}
\vspace{0.5cm}
{\LARGE\bfseries Proposta Técnica e Econômica\par}
\vspace{1cm}
{\Large Fábrica de Tijolos Ecológicos\par}

\vspace{2cm}

\begin{flushleft}
\textbf{Disciplina:} Redes Industriais\\
\textbf{Professor:} Me. Deivison S. Takatu\\
\textbf{Data:} Março de 2026
\end{flushleft}

\vspace{2cm}

\begin{flushright}
\textbf{Aluno:} Felipe Marques\\
\textbf{Curso:} Análise e Desenvolvimento de Sistemas\\
\textbf{Instituição:} SENAI
\end{flushright}

\vfill

\end{titlepage}

% Sumário
\tableofcontents
\newpage

% Conteúdo do Projeto
\section{Introdução}

O presente projeto consiste em uma proposta técnica e econômica completa para a implementação de uma infraestrutura de rede industrial em uma fábrica de tijolos ecológicos. O objetivo é criar um sistema de automação robusto, escalável e integrado, que permita o controle, monitoramento e a futura integração com sistemas corporativos (ERP), seguindo os fundamentos de redes industriais e os requisitos críticos de ambientes produtivos.

A fabricação de tijolos ecológicos é um processo que exige controle preciso de variáveis como temperatura, umidade, pressão e tempo de cura. A automação adequada desses processos resulta em ganhos significativos de qualidade, produtividade e sustentabilidade, além de reduzir desperdícios e consumo de energia.

\section{Descrição do Ambiente Produtivo}

\subsection{Características da Fábrica de Tijolos Ecológicos}

A fábrica de tijolos ecológicos é um ambiente industrial com características específicas que impactam diretamente o projeto de rede:

\begin{itemize}
    \item \textbf{Processo Produtivo:} Mistura de solo-cimento, prensagem, cura a vapor e armazenamento
    \item \textbf{Condições Ambientais:} Alta umidade (câmaras de cura), presença de poeira (solo/brita), vibração (prensa hidráulica)
    \item \textbf{Temperatura:} Variação de 15°C a 85°C (câmaras de cura)
    \item \textbf{Layout:} Galpão industrial de aproximadamente 2.000 m², dividido em 4 áreas principais
\end{itemize}

\subsection{Requisitos Críticos do Ambiente}

Conforme estudado na disciplina, redes industriais devem atender requisitos críticos que diferenciam-no de redes corporativas:

\begin{table}[h]
\centering
\begin{tabular}{@{}p{4cm}|p{10cm}@{}}
\toprule
\textbf{Requisito} & \textbf{Descrição e Aplicação} \\
\midrule
\textbf{Robustez} & Equipamentos devem ser industriais (não adaptados de escritório), com tolerância a ruído elétrico, vibração e poeira. Necessário grau de proteção IP67 para áreas de cura. \\
\midrule
\textbf{Determinismo} & A comunicação deve garantir tempo de resposta previsível para sincronização dos processos de prensagem e cura. \\
\midrule
\textbf{Disponibilidade} & A arquitetura deve priorizar o uptime, evitando pontos únicos de falha que parem a produção. \\
\midrule
\textbf{Segurança} & Proteção contra acesso não autorizado e falhas que podem causar danos materiais ou pessoais. \\
\bottomrule
\end{tabular}
\caption{Requisitos críticos aplicados à fábrica de tijolos}
\end{table}

\section{Plano Orçamentário e Levantamento Técnico}

Esta seção apresenta o levantamento detalhado dos equipamentos necessários para a infraestrutura de comunicação e automação da fábrica de tijolos ecológicos.

\subsection{Controladores Lógicos Programáveis (CLPs)}

Os CLPs constituem o cérebro do sistema de controle, responsáveis pela leitura de sensores e atuação sobre os atuadores.

\begin{table}[h]
\centering
\begin{tabular}{@{}p{3cm}|p{5cm}|p{2cm}|p{3cm}@{}}
\toprule
\textbf{Modelo} & \textbf{Especificações} & \textbf{Qtd} & \textbf{Valor Unitário} \\
\midrule
Siemens S7-1200 & CPU 1214C DC/DC/DC, 14 DI, 10 DO, 2 AI, PROFINET & 3 & R\$ 4.500,00 \\
\midrule
Siemens S7-1200 & Módulo SM 1223, 16 DI/16 DO 24VDC & 4 & R\$ 880,00 \\
\midrule
Siemens S7-1200 & Módulo SM 1231, 4 AI 13-bit & 2 & R\$ 1.200,00 \\
\midrule
Siemens S7-1200 & Cartão Memória 2GB & 3 & R\$ 350,00 \\
\bottomrule
\end{tabular}
\caption{CLPs e módulos de expansão -- Fonte: Siemens Industrial Automation}
\end{table}

\textbf{Justificativa:} A linha S7-1200 foi escolhida por ser uma plataforma compacta e modular, ideal para aplicações de automação de pequeno a médio porte. O protocolo PROFINET integrado permite comunicação determinística via Ethernet, essencial para sincronização dos processos de prensagem.

\subsection{Switches Industriais}

Os switches industriais são responsáveis pela interconexão dos dispositivos na rede de chão de fábrica.

\begin{table}[h]
\centering
\begin{tabular}{@{}p{3cm}|p{5cm}|p{2cm}|p{3cm}@{}}
\toprule
\textbf{Modelo} & \textbf{Especificações} & \textbf{Qtd} & \textbf{Valor Unitário} \\
\midrule
Allen-Bradley & Stratix 2000, 8 portas 10/100, RJ45, IP30 & 6 & R\$ 2.800,00 \\
\midrule
Allen-Bradley & Stratix 2000, 5 portas + fibra, IP67 & 4 & R\$ 4.200,00 \\
\midrule
Cisco/AB & Módulo SFP 100FX, Fibra Multimodo & 8 & R\$ 650,00 \\
\bottomrule
\end{tabular}
\caption{Switches industriais -- Fonte: Rockwell Automation}
\end{table}

\textbf{Justificativa:} A linha Stratix 2000 oferece robustez necessária para ambientes industriais com certificação IP30 (gabinete) e IP67 (campo). A opção com fibra óptica para interligação de áreas permite maior distância e imunidade a interferências eletromagnéticas.

\subsection{Sensores e Transmissores}

Os sensores são os "olhos e ouvidos" do sistema, responsáveis por capturar as variáveis do processo produtivo.

\begin{table}[h]
\centering
\begin{tabular}{@{}p{3cm}|p{5cm}|p{2cm}|p{3cm}@{}}
\toprule
\textbf{Modelo} & \textbf{Especificações} & \textbf{Qtd} & \textbf{Valor Unitário} \\
\midrule
Supmea & Transmissor Pressão SUP-P300, 4-20mA, 0-10 Bar & 8 & R\$ 850,00 \\
\midrule
ACAT & Transmissor Temperatura PT100, -50 a 400°C, IP65 & 12 & R\$ 380,00 \\
\midrule
Weilong & Sensor Nível Ultrassônico, 4-20mA, 0-10m & 6 & R\$ 650,00 \\
\midrule
SICK & Sensor Indutivo 12mm, M12, PNP & 24 & R\$ 220,00 \\
\midrule
SICK & Fotoelétrico Retrorefletivo, M18, 4m & 16 & R\$ 285,00 \\
\bottomrule
\end{tabular}
\caption{Sensores e transmissores -- Fonte: Pesquisa de mercado 2026}
\end{table}

\textbf{Justificativa:} Sensores com saída 4-20mA foram escolhidos por sua imunidade a ruído elétrico e compatibilidade universal com CLPs industriais. A proteção IP65/IP67 garante operação em ambientes úmidos e empoeirados.

\subsection{Interfaces Homem-Máquina (IHMs)}

As IHMs permitem a interação entre operadores e o sistema de controle, proporcionando visualização e parâmetros.

\begin{table}[h]
\centering
\begin{tabular}{@{}p{3cm}|p{5cm}|p{2cm}|p{3cm}@{}}
\toprule
\textbf{Modelo} & \textbf{Especificações} & \textbf{Qtd} & \textbf{Valor Unitário} \\
\midrule
Siemens & SIMATIC HMI TP700, 7" Touch, PROFINET & 3 & R\$ 6.400,00 \\
\midrule
Siemens & SIMATIC HMI KTP400, 4" Touch, PROFINET & 4 & R\$ 3.200,00 \\
\midrule
Weinview & MT8051iP, 4.3", Ethernet, IP65 & 2 & R\$ 1.350,00 \\
\bottomrule
\end{tabular}
\caption{Interfaces Homem-Máquina -- Fonte: Siemens/Weinview}
\end{table}

\textbf{Justificativa:} As IHMs da linha Comfort da Siemens oferecem alta durabilidade e integração nativa com CLPs S7-1200 via PROFINET. Os modelos de 7" serão instalados em painéis de controle principais, enquanto os de 4" em pontos locais de operação.

\subsection{Sistema SCADA/MES}

O sistema SCADA (Supervisory Control and Data Acquisition) é responsável pela supervisão centralizada e coleta de dados para análise histórica.

\begin{table}[h]
\centering
\begin{tabular}{@{}p{3cm}|p{5cm}|p{2cm}|p{3cm}@{}}
\toprule
\textbf{Software} & \textbf{Especificações} & \textbf{Qtd} & \textbf{Valor Total} \\
\midrule
AVEVA Edge & Licença Desenvolvimento + Runtime, 1500 tags & 1 & R\$ 18.500,00 \\
\midrule
AVEVA Edge & Thin Client, 5 usuários & 1 & R\$ 4.200,00 \\
\midrule
InPlant SCADA & Web Publisher, módulo adicional & 1 & R\$ 6.800,00 \\
\bottomrule
\end{tabular}
\caption{Software SCADA -- Fonte: AVEVA/Supcon}
\end{table}

\textbf{Justificativa:} O AVEVA Edge (anteriormente InduSoft Web Studio) é uma plataforma HMI/SCADA flexível e escalável. A opção por Thin Client permite acesso via navegador web, facilitando a supervisão remota. O módulo Web Publisher possibilita a visualização de telas via internet/intranet.

\subsection{Infraestrutura de Rede e Cabeamento}

\begin{table}[h]
\centering
\begin{tabular}{@{}p{3cm}|p{5cm}|p{2cm}|p{3cm}@{}}
\toprule
\textbf{Item} & \textbf{Especificações} & \textbf{Qtd} & \textbf{Valor Unitário} \\
\midrule
Cabo Ethernet & Industrial Cat6, LSZH, blindado, 305m & 4 & R\$ 1.850,00 \\
\midrule
Cabo Fibra & Multimodo OM3, 4 fibras, 1.000m & 1 & R\$ 4.200,00 \\
\midrule
Conectores & RJ45 Industrial, blindado, metal & 100 & R\$ 28,00 \\
\midrule
Racks & Paredes 19", 6U, IP65, industrial & 4 & R\$ 890,00 \\
\midrule
Organizadores & Guia de cabos, metálico, 2m & 20 & R\$ 65,00 \\
\bottomrule
\end{tabular}
\caption{Infraestrutura de rede -- Fonte: Cotação mercado}
\end{table}

\subsection{Resumo do Orçamento}

\begin{table}[h]
\centering
\begin{tabular}{@{}p{6cm}|p{3cm}|p{4cm}@{}}
\toprule
\textbf{Categoria} & \textbf{Valor Total} & \textbf{\% do Total} \\
\midrule
CLPs e Módulos & R\$ 18.940,00 & 27,1\% \\
\midrule
Switches Industriais & R\$ 31.000,00 & 44,3\% \\
\midrule
Sensores e Transmissores & R\$ 15.580,00 & 22,3\% \\
\midrule
Interfaces HMI & R\$ 25.400,00 & 36,3\% \\
\midrule
Software SCADA & R\$ 29.500,00 & 42,2\% \\
\midrule
Infraestrutura de Rede & R\$ 9.430,00 & 13,5\% \\
\midrule
\textbf{TOTAL DO PROJETO} & \textbf{R\$ 69.850,00} & \textbf{100\%} \\
\bottomrule
\end{tabular}
\caption{Resumo do orçamento total do projeto}
\end{table}

\section{Análise de Adequação e Escalabilidade}

\subsection{Adequação Ambiental}

Todos os equipamentos selecionados possuem certificações adequadas para o ambiente da fábrica de tijolos:

\begin{table}[h]
\centering
\begin{tabular}{@{}p{4cm}|p{3cm}|p{3cm}|p{5cm}@{}}
\toprule
\textbf{Equipamento} & \textbf{Proteção IP} & \textbf{Temperatura} & \textbf{Observações} \\
\midrule
CLPs Siemens S7-1200 & IP20 & -20 a 60°C & Instalados em painéis \\
\midrule
Switches Stratix & IP30/IP67 & -40 a 75°C & Campo: IP67 \\
\midrule
Sensores & IP65/IP67 & -40 a 85°C & Adequados câmara cura \\
\midrule
IHMs & IP65/IP66 & 0 a 50°C & Painel selado \\
\bottomrule
\end{tabular}
\caption{Adequação ambiental dos equipamentos}
\end{table}

\textbf{Proteção contra Interferência Elétrica:}
\begin{itemize}
    \item Cabos blindados para todos os sinais analógicos e digitais
    \item Switches com portas RJ45 industriais com proteção ESD
    \item Isolação galvânica em todos os transmissores 4-20mA
    \item Aterramento conforme norma IEC 60364
\end{itemize}

\subsection{Escalabilidade do Sistema}

A arquitetura proposta permite expansão futura com mínimo impacto:

\begin{itemize}
    \item \textbf{Rede PROFINET:} Permite adição de até 64 dispositivos por segmento
    \item \textbf{CLPs:} Capacidade para até 8 módulos de expansão por CPU
    \item \textbf{SCADA:} Licença expansível até 64.000 tags
    \item \textbf{Switches:} Portas SFP disponíveis para expansão via fibra
\end{itemize}

\textbf{Cenários de Expansão:}
\begin{enumerate}
    \item Adição de nova linha de prensagem: +1 CLP, +6 sensores, +1 IHM
    \item Ampliação da área de cura: +12 sensores T/U, +1 CLP remoto
    \item Integração com ERP: Adição de gateway OPC UA sem modificar CLPs
\end{enumerate}

\subsection{Compatibilidade Corporativa}

A solução proposta permite integração vertical futura com sistemas de gestão:

\begin{itemize}
    \item \textbf{OPC UA:} Todos os CLPs S7-1200 suportam OPC UA nativo
    \item \textbf{MQTT:} SCADA com publicação para sistemas de nuvem
    \item \textbf{SQL:} Histórico de dados diretamente em banco de dados corporativo
    \item \textbf{REST API:} Interface para integração com sistemas web
\end{itemize}

\section{Análise de Viabilidade Técnica e Econômica}

\subsection{Análise Custo-Benefício}

\begin{table}[h]
\centering
\begin{tabular}{@{}p{5cm}|p{5cm}|p{4cm}@{}}
\toprule
\textbf{Benefício} & \textbf{Descrição} & \textbf{Impacto Estimado} \\
\midrule
Redução de Refugo & Monitoramento preciso dos parâmetros de cura & -15\% desperdício \\
\midrule
Aumento Produtividade & Automação dos ciclos de prensagem & +25\% produção \\
\midrule
Economia Energia & Otimização câmaras de cura & -20\% consumo energia \\
\midrule
Melhoria Qualidade & Padronização dos processos & +10\% qualidade \\
\midrule
Redução Paradas & Manutenção preditiva baseada em dados & -30\% paradas \\
\bottomrule
\end{tabular}
\caption{Benefícios esperados da implementação}
\end{table}

\subsection{Cálculo do Retorno sobre Investimento (ROI)}

\textbf{Cenário Base (sem automação):}
\begin{itemize}
    \item Produção mensal: 50.000 tijolos
    \item Preço unitário: R\$ 1,20
    \item Custo produção: R\$ 0,90/unid
    \item Refugo atual: 8\%
    \item Lucro mensal: R\$ 10.800,00
\end{itemize}

\textbf{Cenário com Automação:}
\begin{itemize}
    \item Produção mensal: 62.500 tijolos (+25\%)
    \item Refugo reduzido para: 3\%
    \item Economia energia: R\$ 800,00/mês
    \item Lucro mensal projetado: R\$ 20.325,00
\end{itemize}

\textbf{Cálculo do ROI:}
\begin{itemize}
    \item Investimento total: R\$ 69.850,00
    \item Aumento mensal de lucro: R\$ 9.525,00
    \item \textbf{Payback:} 7,3 meses
    \item \textbf{ROI anual:} 163,7\%
\end{itemize}

\subsection{Gestão de Riscos}

\begin{table}[h]
\centering
\begin{tabular}{@{}p{4cm}|p{4cm}|p{4cm}|p{3cm}@{}}
\toprule
\textbf{Risco} & \textbf{Probabilidade} & \textbf{Impacto} & \textbf{Mitigação} \\
\midrule
Falha CLP principal & Baixa & Alto & Redundância software \\
\midrule
Interferência rede & Média & Médio & Cabos blindados \\
\midrule
Obsolescência software & Baixa & Alto & Escolha fornecedor estável \\
\midrule
Dependência fornecedor & Média & Médio & Múltiplos fornecedores \\
\midrule
Falha treinamento & Alta & Médio & Treinamento inicial \\
\bottomrule
\end{tabular}
\caption{Matriz de riscos do projeto}
\end{table}

\subsection{Análise de Sensibilidade}

Considerando variações no cenário econômico:

\begin{table}[h]
\centering
\begin{tabular}{@{}p{4cm}|p{3cm}|p{3cm}|p{3cm}|p{3cm}@{}}
\toprule
\textbf{Cenário} & \textbf{Produtividade} & \textbf{Payback} & \textbf{ROI (12 meses)} \\
\midrule
Otimista & +35\% & 5,8 meses & 207\% \\
\midrule
Base (Projetado) & +25\% & 7,3 meses & 164\% \\
\midrule
Conservador & +15\% & 10,5 meses & 114\% \\
\bottomrule
\end{tabular}
\caption{Análise de sensibilidade -- variações de produtividade}
\end{table}

\section{Arquitetura da Rede Proposta}

\subsection{Topologia de Rede}

A arquitetura proposta segue o modelo de três níveis da pirâmide de automação, interconectados via Ethernet Industrial:

\begin{table}[h]
\centering
\begin{tabular}{@{}p{3cm}|p{4cm}|p{3cm}|p{5cm}@{}}
\toprule
\textbf{Nível} & \textbf{Protocolo} & \textbf{Equipamentos} & \textbf{Função} \\
\midrule
Nível 2 (SCADA) & TCP/IP & Servidor SCADA, clientes & Supervisão \\
\midrule
Nível 1 (Controle) & PROFINET & CLPs, IHMs, Switches & Controle lógico \\
\midrule
Nível 0 (Campo) & 4-20mA/Digital & Sensores, Atuadores & Aquisição dados \\
\bottomrule
\end{tabular}
\caption{Níveis da arquitetura de rede}
\end{table}

\subsection{Diagrama de Arquitetura}

\textbf{Descrição da Topologia:}

\begin{enumerate}
    \item \textbf{Nível Corporativo (Topo):}
    \begin{itemize}
        \item Servidor SCADA/Historiador
        \item Estações de trabalho do gerenciamento
        \item Conexão com ERP (futuro)
    \end{itemize}

    \item \textbf{Nível de Controle (Meio):}
    \begin{itemize}
        \item 3 CLPs Siemens S7-1200 controlando áreas:
        \begin{itemize}
            \item Área 1: Mistura e dosagem
            \item Área 2: Prensagem
            \item Área 3: Cura e estocagem
        \end{itemize}
        \item 7 IHMs distribuídas nos pontos de operação
        \item Switches industriais formando backbone anel
    \end{itemize}

    \item \textbf{Nível de Campo (Base):}
    \begin{itemize}
        \item 26 sensores de temperatura/umidade
        \item 8 transmissores de pressão
        \item 6 sensores de nível
        \item 40 sensores discretos (indutivos/fotoelétricos)
    \end{itemize}
\end{enumerate}

\subsection{Segurança da Rede}

A segurança da rede industrial é implementada em camadas:

\begin{table}[h]
\centering
\begin{tabular}{@{}p{4cm}|p{9cm}@{}}
\toprule
\textbf{Camada} & \textbf{Medidas de Segurança} \\
\midrule
\textbf{Física} & Gabinetes trancados, controle de acesso, câmeras \\
\midrule
\textbf{Rede} & VLANs separadas (chão de fábrica x corporativo), firewall industrial \\
\midrule
\textbf{Aplicação} & Autenticação por usuário, níveis de acesso, logs de auditoria \\
\midrule
\textbf{Dados} & Comunicação criptografada, backup diário do histórico \\
\bottomrule
\end{tabular}
\caption{Medidas de segurança por camada}
\end{table}

\section{Conclusão}

 Com um investimento total de R\$ 69.850,00, o projeto apresenta um payback estimado de 7,3 meses e ROI de 164\% no primeiro ano, demonstrando atratividade econômica mesmo em cenários conservadores.

\subsection{Benefícios Esperados}

A implementação do sistema proposto resultará em:

\begin{itemize}
    \item \textbf{Ganho de Produtividade:} Aumento de 25\% na capacidade produtiva através de automação dos ciclos
    \item \textbf{Redução de Custos:} Diminuição de 20\% no consumo energético e 15% no índice de refugo
    \item \textbf{Melhoria de Qualidade:} Padronização dos processos produtivos com controle estatístico
    \item \textbf{Base para Crescimento:} Arquitetura escalável que permite futuras expansões
\end{itemize}

\subsection{Retorno sobre Investimento}

Com um investimento total de R\$ 69.850,00, o projeto apresenta um payback estimado de 7,3 meses e ROI de 164\% no primeiro ano, demonstrando atratividade econômica mesmo em cenários conservadores.

\subsection{Alinhamento com Conceitos da Disciplina}

O projeto aplica corretamente os conceitos de redes industriais estudados:

\begin{itemize}
    \item \textbf{Integração Vertical:} Do sensor (Nível 0) ao SCADA (Nível 2), permitindo tomada de decisão baseada em dados em tempo real
    \item \textbf{Requisitos Críticos:} Robustez, determinismo e disponibilidade atendidos através de equipamentos industriais certificados
    \item \textbf{Segregação de Redes:} Separação entre rede industrial e rede corporativa, preparada para futura integração via OPC UA
\end{itemize}

A fábrica de tijolos ecológicos, ao implementar este projeto, estará posicionada para competir efetivamente no mercado através da moderna automação industrial, alinhada aos conceitos da Indústria 4.0 e preparada para integrações futuras com sistemas de gestão corporativa.

\section*{Referências Bibliográficas}

\begin{enumerate}
    \item \textbf{SIEMENS.} SIMATIC S7-1200 System Manual. Disponível em: \url{https://support.industry.siemens.com/}. Acesso em: mar. 2026.

    \item \textbf{ROCKWELL AUTOMATION.} Stratix 2000 Unmanaged Switches User Manual. Disponível em: \url{https://www.rockwellautomation.com/}. Acesso em: mar. 2026.

    \item \textbf{AVEVA.} AVEVA Edge Technical Documentation. Disponível em: \url{https://www.aveva.com/}. Acesso em: mar. 2026.

    \item \textbf{TAKATU, Deivison S.} Redes Industriais -- Fundamentos e Aplicações. Material didático, 2024.

    \item \textbf{IEC 60529.} Degrees of Protection Provided by Enclosures (IP Code). International Electrotechnical Commission, 1989.

    \item \textbf{IEC 61158.} Industrial Communication Networks -- Fieldbus Specifications. International Electrotechnical Commission, 2019.

    \item \textbf{LARA, Carla Eduarda Orlando de Moraes de.} Automação e controle industrial. Curitiba: Contentus, 2021.

    \item \textbf{KAGERMANN, Henning et al.} Recommendations for Implementing Strategic Initiative Industrie 4.0. Final Report of the Industrie 4.0 Working Group, 2013.

    \item \textbf{PI (PROFIBUS \& PROFINET INTERNATIONAL).} PROFINET Technology and Application. System Manual, 2020.
\end{enumerate}

\end{document}
